\documentclass[10pt,a4paper]{article}

\usepackage[utf8]{inputenc}
\usepackage[T1]{fontenc}
\usepackage[portuguese]{babel}
\usepackage[margin=1.35cm]{geometry}
\usepackage[hidelinks]{hyperref}
\usepackage{enumitem}
\usepackage{titlesec}

\setlist[itemize]{leftmargin=1.1em,itemsep=1pt,topsep=1pt}
\setlength{\parindent}{0pt}
\setlength{\parskip}{2pt}
\titlespacing*{\section}{0pt}{6pt}{3pt}

\begin{document}

{\Large\textbf{Renan Rodrigues}}\\
Engenharia da Computação \textbar{} Full Stack Developer

\vspace{2pt}
São Paulo, SP - Brasil \textbar{}
Telefone: (11) 93340-7053 \textbar{}
Email: \href{mailto:renanrdemeneses@gmail.com}{renanrdemeneses@gmail.com}

\href{https://renanrod.vercel.app/}{Portfolio} \textbar{}
\href{https://linkedin.com/in/renanrod4}{LinkedIn} \textbar{}
\href{https://github.com/renanrod4}{GitHub}

\vspace{2pt}\hrule\vspace{2pt}

\section*{Resumo}
Estudante de Engenharia da Computação focado em desenvolvimento Full Stack. Experiência na criação de aplicações web, APIs REST e sistemas interativos utilizando JavaScript/TypeScript, React, Next.js e Rust. Busco aplicar meus conhecimentos em projetos reais e evoluir continuamente como desenvolvedor.

\section*{Formação}
\textbf{Bacharelado em Engenharia de Computação} \hfill 2025 -- 2030\\
Universidade de Sorocaba (UNISO)

\textbf{Ensino Médio Completo} \hfill \textbf{2022 -- 2024}\\
E.E. Bel Elias Alves da Costa Bacharel\hfill 2023 -- 2024\\
Colégio Cognos Sistema Mackenzie de Ensino\hfill{2022}\\

\section*{Experiência Acadêmica e P\&D (UNISO)}
\textbf{Uniso Flow} \textbar{} Full Stack Developer \hfill \textbf{2025 -- Presente}\\
Sistema de análise preditiva para gestão de fluxo acadêmico, com cruzamento de pré-requisitos para prever demanda e identificar gargalos.\\
\textbf{Tech:} Next.js, Tailwind, Radix UI, Tauri (Rust)

\textbf{Termômetro Laser Industrial} \textbar{} ESP32 \& C++ \hfill \textbf{2025 -- Presente}\\
Prototipagem de hardware para medição térmica de precisão (vacinas e alimentos), com firmware para calibração de sensores ópticos e processamento de sinais.\\
\textbf{Tech:} C++, ESP32, Sensores Infravermelhos


\section*{Projetos Técnicos}
\textbf{\href{https://octodev.vercel.app/}{OctoDev}} \textbar{} Tauri, Rust, React.js, Next.js, MongoDB, JSON\\
Plataforma gamificada de aprendizado com desafios interativos e árvores de habilidades.

\textbf{\href{https://flexy-api.vercel.app}{FlexyAPI}} \textbar{} React.js, Next.js, REST APIs\\
Uma Mock Data API orientada por IA. Gera estruturas JSON complexas e contextuais para prototipagem rápida e testes de frontend.

\textbf{\href{https://rr-dealership.vercel.app/}{RR's Dealership}} \textbar{} React.js, Next.js, JSON\\
Catálogo dinâmico com processamento de dados estruturados para mais de 70 veículos.

\textbf{\href{https://github.com/renanrod4/Periodic-Table-game}{Periodic Table Game}} \textbar{} Python, Pygame\\
Quiz interativo e gamificado para memorização dos elementos da tabela periódica.

\section*{Skills}
\textbf{Software e Ferramentas:} Blender, Unity, VS Code, Arduino, Figma, Git, GitHub,\\
Bootstrap, Tailwind, Vercel

\textbf{Desenvolvimento Web:} HTML, CSS, JavaScript, TypeScript, React, Next.js, Node.js, PHP, REST

\textbf{Linguagens:} Python, Rust, C++, C\#, Bash

\textbf{Bancos de Dados:} MySQL, MongoDB

\section*{Idiomas}
Português (Nativo), Inglês (Avançado), Alemão (Básico)

\section*{Conquistas}
\begin{itemize}
  \item Feira de Ciências de SP (2023): Pesquisa de múons (Fermilab/USP).
  \item Medalha de Bronze: Olimpíada de Matemática (2024).
  \item Medalhista: Olimpíada Brasileira de Astronomia (2024).
\end{itemize}

\end{document}
